% ===== 原文 pp.16–18 =====
% tr-03 块，第 1 片：
%   §3 A class of four-dimensional theories（标题）
%   §3.1 N=2 supersymmetric gauge theories
%   §3.2 Class S theories of type SU(2)（标题）+ §3.2.1 + §3.2.2
%   图 7（p.17）、图 8（p.18）；tag 3.1–3.6 中的 3.1–3.6 全在本片
% 注意：p.18 末尾那句 “We can give a mass term using SU(2) flavor symmetry,
%       still preserving N=2 supersymmetry.” 跨页到 p.19，为不切断句子，
%       整句移到 tr03b.tex 开头。

% [原文 p.16]

\section{一类四维理论}

本节要引入一类四维 $\cN=2$ 超对称规范理论，文献中通常称之为 $\SU(2)$ 型的 class S 理论。这一类理论最早由 [4] 引入；详尽的综述可参见 \eg\ [20]。我们先快速回顾一下 $\cN=2$ 拉格朗日量最基本的内容。

% [Tutor Note: p.16，§3 开头。“class S theories of type SU(2)” 里的
%  “type” 说的是规范群取 SU(2) 这一型（就 6d (2,0) 理论的分类而言是
%  A_1 型），不是另立一族的理论。原文紧接着的 [4] 即 Gaiotto,
%  “N=2 Dualities”（JHEP 1208 (2012) 034, arXiv:0904.2715），本节的
%  整套几何构造（图 7–图 10 与 (3.3)–(3.11)）都出自该文；[20] 是
%  Tachikawa 自己的讲义 “N=2 supersymmetric dynamics for pedestrians”。]

\subsection{$\cN=2$ 超对称规范理论}

我们假定读者了解 $\cN=1$ 超场的基本知识。我们所处理的 $\cN=2$ 理论可以这样得到：在 $\cN=1$ 形式中显然可见的超对称之外，再施加一个与它\emph{不}对易的 $\SU(2)_R$ 对称性。

让我们从 $\cN=2$ 矢量多重态讲起。它由规范群 $G$ 的一个 $\cN=1$ 矢量多重态 $V$，连同 $G$ 的伴随表示中的一个 $\cN=1$ 手征多重态 $\Phi$ 组成。我们考虑拉格朗日量
\begin{equation}
  \frac{\Im\tau}{4\pi}\int \dd^4\theta\, \tr \Phi^\dagger \ee^{[V,\cdot]}\Phi
  + \int \dd^2\theta\, \frac{-\ii}{8\pi}\tau\, \tr W_\alpha W^\alpha + cc.
  \tag{3.1}
\end{equation}
其中 $\tau = 4\pi\ii/g^2 + \theta/2\pi$ 是复化规范耦合。把超场展开成分量便能看到，$V$ 中的规范微子（gaugino）$\lambda$ 与 $\Phi$ 中的手征微子（chiralino）$\psi$ 与其余场具有完全相同的耦合，于是 $\SU(2)_R$ 对称性得以实现。

% [Tutor Note: p.16，式 (3.1)。$\tau = 4\pi i/g^2+\theta/2\pi$ 的实部是
%  $\theta/2\pi$、虚部是 $4\pi/g^2$，所以 $\Im\tau/4\pi = 1/g^2$：这正是
%  第一项（$d^4\theta$ 的 Kähler 势项）的系数，也就是通常的 $1/g^2$。
%  第二项的系数若把 $\tau$ 拆开则是
%  $-i\tau/8\pi = \dfrac{1}{2g^2}-\dfrac{i\theta}{16\pi^2}$，
%  前一半是规范动能、后一半给出 $\theta$ 角项；因此 $\theta$ 与 $1/g^2$
%  的相对系数完全由同一个复化耦合 $\tau$ 决定。抄写时最易出错的是
%  把第一项的 $\Im\tau/4\pi$ 误抄成 $\tau/4\pi$（丢掉 $\Im$），
%  或把两个分母 $4\pi$、$8\pi$ 弄反。]

接下来引入 $\cN=2$ 超多重态，它处于规范群 $G$ 的表示 $R$ 之中。它由分别处于表示 $R$ 与 $\bar R$ 的一对 $\cN=1$ 手征多重态 $Q$、$\tilde Q$ 组成。其拉格朗日量为
\begin{equation}
  \int \dd^4\theta\,\bigl(Q^\dagger \ee^{V} Q + \tilde Q\, \ee^{-V}\tilde Q^\dagger\bigr)
  + \Bigl(\int \dd^2\theta\, \tilde Q\Phi Q + cc.\Bigr)
  \tag{3.2}
\end{equation}
其中 $\mu$ 是质量项。这里，$\SU(2)_R$ 对称性转动的是 $Q$ 与 $\tilde Q^\dagger$ 的标量分量。

% [Tutor Note: p.16，式 (3.2) 下方 “where $\mu$ is the mass term” 一句。
%  已按 900 dpi 复核原文：式 (3.2) 的通篇（两处积分）里并没有 $\mu$，
%  整个 p.16 上 $\mu$ 也只在“这一句”里出现一次，而紧接着的 p.17 开头
%  又明说“要给出质量项，只需给（下划线的）$\underline{\Phi}$ 一个真空
%  期望值”。故这一句应是作者从早期版本（当时式 (3.2) 里带有
%  $+\int d^2\theta\,\mu\tilde QQ$ 这样的项）残留下来的一句；本译本正文
%  照原文排，读时把 “$\mu$ 是质量项” 当作无效说明即可。]

% [原文 p.17]

上面的拉格朗日量描述的是无质量超多重态。要给出质量项，我们可以让上面拉格朗日量中的 $\Phi$ 取一个真空期望值。例如取一对超多重态 $Q^a_i$ 与 $\tilde Q^i_a$，其中 $a=1,\dots,N_c$、$i=1,\dots,N_f$。它们处于 $\SU(N_c)\times\U(N_f)$ 的双基表示，因此在拉格朗日量中有如下耦合
\begin{equation}
  \tilde Q^i_a\, \Phi^a{}_b\, Q^b_i
  + \tilde Q^i_a\, \underline{\Phi}^j{}_i\, Q^a_j
  \tag{3.3}
\end{equation}
其中 $\Phi$ 处于 $\SU(N_c)$ 的伴随表示，$\underline{\Phi}$ 处于 $\U(N_f)$ 的伴随表示。现在我们就把 $\underline{\Phi}$ 及其相应的 $\cN=1$ 矢量多重态当作外部的背景场，直接赋予真空期望值 $\langle\underline{\Phi}\rangle^i_j = m^i_j$。于是我们得到形如 $m^i{}_j\tilde Q^i Q_j$ 的质量项。为保持 $\SU(2)_R$ 不变性，我们要求 $[m,m^\dagger]=0$，也就是说 $m$ 可对角化。已知这是给超多重态赋予质量的唯一办法。

% [Tutor Note: p.17 第一行 “we can give a vev to Φ in the Lagrangian above”
%  里的那个 $\Phi$ 原文没有加下划线（已按 1600 dpi 逐字复核），而同一段
%  后面两处 $\langle\underline{\Phi}\rangle^i_j=m^i_j$ 与 “$\underline{\Phi}$
%  and its associated N=1 vector multiplet” 都加了。按式 (3.3)，能给
%  $(i,j)$ 指标提供质量项的只能是 $\U(N_f)$ 伴随的那个
%  $\underline{\Phi}$，故第一处疑为漏排下划线。译文照原文排（不加下划线），
%  读者按上下文理解为 $\underline{\Phi}$ 即可。]

% [Tutor Note: p.17，式 (3.3)。这里出现的两个 $\Phi$ 不是同一个场，原文
%  用下划线区分：$\Phi^a{}_b$ 是 $\SU(N_c)$ 的伴随手征场（指标 $a,b$），
%  $\underline{\Phi}^j{}_i$ 是 $\U(N_f)$ 的伴随手征场（指标 $i,j$）。
%  抄写时最易漏掉第二处的下划线。指标配对可逐个核对：第一项
%  $\tilde Q^i_a\,\Phi^a{}_b\,Q^b_i$ 中 $\tilde Q$ 的下 $a$ 与 $\Phi$ 的上 $a$
%  缩并、$\Phi$ 的下 $b$ 与 $Q$ 的上 $b$ 缩并、$\tilde Q$ 的上 $i$ 与 $Q$ 的下
%  $i$ 缩并；第二项则是 $i$ 走 $\underline{\Phi}$、$a$ 走 $\tilde Q$ 与 $Q$ 之间。]

% [Tutor Note: p.17，质量项的来历（原文省略的代数）。把
%  $\langle\underline{\Phi}\rangle^j{}_i=m^j{}_i$ 代入式 (3.3) 第二项，
%  得 $\tilde Q^i_a m^j{}_i Q^a_j = m^j{}_i(\tilde Q^i_a Q^a_j)$；把括号里
%  缩并掉 $a$ 之后的量记为 $\tilde Q^iQ_j$ 并重新命名指标，即得正文所说的
%  $m^i{}_j\tilde Q^iQ_j$。再要求 $\SU(2)_R$ 不变，就需要 $m$ 与 $m^\dagger$
%  对易（否则 $\SU(2)_R$ 会把这个质量项变成另一个不可对角的组合），
%  故 $m$ 可对角化、可写成 $\diag(m_1,\dots,m_{N_f})$。]

当 $m=0$ 时，$\U(N_f)$ 对称性是一个整体的味对称性。对一般的非零对角 $m$，这一 $\U(N_f)$ 对称性被进一步明显破缺为 $\U(1)^{N_f}$。我们通常说，这个质量项 $m^i{}_j\tilde Q^i Q_j$ 是与味对称性 $\U(N_f)$ 相联系的。我们也常滥用术语，即使 $m\neq 0$ 也说该理论具有 $\U(N_f)$ 味对称性。

最后引入半超多重态。当表示 $R$ 是伪实的，\ie 当 $R$ 与 $\bar R$ 作为表示等价、且存在规范不变的反对称二形式 $\epsilon_{ab}$ 时，我们可以施加条件
\begin{equation}
  Q_a = \epsilon_{ab}(\tilde Q^\dagger)^b
  \tag{3.4}
\end{equation}
它与 $\SU(2)_R$ 不变性相容。这称为处于表示 $R$ 中的半超多重态。注意它只由处于表示 $R$ 中的单个 $\cN=1$ 手征多重态组成。

% [Tutor Note: p.17，式 (3.4)。$\epsilon_{ab}$ 必须是规范不变的反对称
%  二形式：对伪实表示，$\bar R\cong R$，故 $(\tilde Q^\dagger)^b$ 与 $Q_a$
%  同属一个表示，才谈得上用 $\epsilon_{ab}$ 把上标降成下标。若 $R$ 是实表示
%  （如伴随表示），也存在不变对称形式但没有不变反对称形式，此时条件
%  (3.4) 不成立，只能有整超多重态；这正是 §3.2.2 里三重态超多重态要用
%  $A_I$ 与 $\tilde A_I$ 两个 $\cN=1$ 手征场来拼的原因。]

\subsection{$\SU(2)$ 型的 class S 理论}

\subsubsection{由三穿刺球面与圆柱面构造}

在第 2 节中，我们回顾了两个二维场论的性质。非常抽象地说，二维场论给每个二维曲面配上一个复数。那里的要点在于求出与三穿刺球面（或三孔球面）相应的振幅，以及与圆柱面相应的振幅。

这里，我们将引入一种办法，给二维曲面配上四维场论而不是复数。同样，要紧的是考虑给三穿刺球面或圆柱面配上什么。

\begin{center}
\begin{minipage}{\textwidth}
\centering
\includegraphics[width=0.45\textwidth]{fig7.png}\par\vspace{2pt}
\figtitle{7}{三穿刺球面对应于 $Q_{aiu}$，具有 $\SU(2)^3$ 对称性。}
\end{minipage}
\end{center}

% [原文 p.18]

首先取一个三穿刺球面，见图 7。我们为三个穿刺点各配一个 $\SU(2)$ 味对称性。$\SU(2)$ 的基础表示是伪实的；而且，奇数个伪实表示的张量积仍是伪实的。因此，处于 $\SU(2)_1\times\SU(2)_2\times\SU(2)_3$ 的表示 $(2,2,2)$ 中的 $\cN=1$ 手征多重态——我们记作
\begin{equation}
  Q_{aiu},\qquad a=1,2;\ i=1,2;\ u=1,2
  \tag{3.5}
\end{equation}
——构成一个半超多重态。这里 $a,i,u$ 分别是 $\SU(2)_{1,2,3}$ 的指标。

% [Tutor Note: p.18，式 (3.5) 之上 “奇数个伪实表示的张量积仍是伪实的”
%  这一步。判据其实就是不变反对称二形式 (3.4) 的存在性：$Q_{aiu}$ 上有三个
%  指标，取 $\epsilon_{a_1a_2}\epsilon_{i_1i_2}\epsilon_{u_1u_2}$ 的张量积，它是
%  $\SU(2)^3$ 不变、并且对整套指标的交换反对称的，于是正好提供了一个
%  规范不变的反对称二形式，使 (3.4) 可以实现。偶数个因子时该张量积对
%  交换对称，给不出 (3.4)，所以 “奇数个” 这个条件不能改成 “任意个”。]

其次取一个圆柱面，见图 8。我们给它配一个复数 $\tau = 4\pi\ii/g^2 + \theta/2\pi$，其中 $g$ 为实数、$\theta \sim \theta + 2\pi$。然后给它配一个规范群为 $\SU(2)$ 的 $\cN=2$ 矢量多重态，其复化耦合为 $\tau$。

\begin{center}
\begin{minipage}{\textwidth}
\centering
\includegraphics[width=0.45\textwidth]{fig8.png}\par\vspace{2pt}
\figtitle{8}{参数为 $\tau$ 的圆柱面对应于一个 $\SU(2)$ 矢量多重态。}
\end{minipage}
\end{center}

现在，给定若干三穿刺球面，我们挑出两个穿刺点。设第一个穿刺点在局部坐标 $z$ 下位于 $z=0$，第二个在局部坐标 $w$ 下位于 $w=0$。我们取一个参数为 $\tau$ 的圆柱面，并作如下认同
\begin{equation}
  zw = \ee^{2\pi\ii\tau}.
  \tag{3.6}
\end{equation}
我们也常用记号 $q = \ee^{2\pi\ii\tau}$。

% [Tutor Note: p.18，式 (3.6) 与 “$q=\ee^{2\pi i\tau}$”。上面的 $\theta$
%  周期条件 $\theta\sim\theta+2\pi$ 等价于 $\tau\sim\tau+1$，而
%  $q=\ee^{2\pi i\tau}$ 在 $\tau\to\tau+1$ 下恰不变——这正是 $q$ 比 $\tau$
%  更适合做模参数的原因（$\tau$ 的 $\tau\to\tau+1$ 是同一圆柱面的大规范
%  变换，不是新曲面）。同理，后面 p.19 的 S-对偶 $\tau\to-1/\tau$ 在 $q$ 上
%  是 $q\to\ee^{-2\pi i/\tau}$，并不是简单的 $q\to1/q$，读 p.20 时不要混。]

相应地，在四维规范理论这一侧，我们作如下操作。由于挑出了两个穿刺点，我们便挑出了两个 $\SU(2)$ 味对称性。现在把它们与一个动力学 $\SU(2)$ 规范场耦合起来，其复化耦合常数为 $\tau$。

\subsubsection{单穿刺环面}

取一个三穿刺球面，把三个穿刺点中的两个用一条参数为 $\tau$ 的圆柱面连接起来。结果是一个带穿刺点、模参数为 $\tau$ 的环面，见图 9。作为规范理论上的操作，我们从三基表示 $Q_{aiu}$ 出发，挑出指标 $a$ 与 $i$，把它与单个 $\SU(2)$ 规范场耦合。现在我们把这作用在 $a$ 与 $i$ 上的两个 $\SU(2)$ 对称性看作同一个；由于二重态乘二重态等于三重态加单态，我们把 $Q_{aiu}$ 重新标号为 $A_{Iu}$（$I=1,2,3$）与 $H_u$，其中指标 $u=1,2$ 仍然对应一个 $\SU(2)$ 味对称性。

我们最终得到的是一个 $\cN=2$ 的 $\SU(2)$ 规范理论，其中含有一个由 $A_I := A_{I,u=1}$、$\tilde A_I := A_{I,u=2}$ 构成的三重态超多重态，以及一个由 $H := H_{u=1}$、$\tilde H := H_{u=2}$ 构成的解耦超多重态。规范群为 $G$、且带有 $G$ 伴随表示超多重态的 $\cN=2$ 规范理论具有增强的 $\cN=4$ 超对称性，通常就径直称作 $\cN=4$ 超杨–米尔斯理论。
